\subsection*{一步一步建立直观}

\paragraph{1. 规则先告诉你下一瞬间往哪里走}

微分方程给出当前位置对应的速度，而不是直接给出整条路线。右图的小箭头标示当地速度方向；选定一个初始点后，轨线必须处处沿着这些方向前进。先选择纯旋转场，比较沿轨线的切向方向与经过位置上的速度箭头是否一致。



\paragraph{2. 同一旋转可以叠加径向收缩}

加入负的径向增长率，轨线向内盘旋。角度继续变化，半径却按指数缩小，因此“在转动”和“向平衡点靠近”并不冲突。这里使用有显式解的线性系统，便于把局部方程、整体轨迹和实时坐标三者直接核对。



\paragraph{3. 初值与稳定性是不同问题}

把径向增长率改为正，同样的旋转伴随向外增长。拖动初始点可得到同一方程的另一条解；修改增长率则已经改变方程本身。靠近原点的初始状态能否始终保持接近，是稳定性问题，不能仅凭某个初值恰好静止来判断。



\subsection*{定义、命题与证明}

\paragraph{定义：线性平面速度场}

考虑 $x'=\lambda x-\omega y$、$y'=\omega x+\lambda y$，其中 $\lambda,\omega$ 为常数。前者控制径向伸缩，后者控制角向转动。

\paragraph{命题：长度平方的增长只由径向项控制}

沿任意可微解，$\dfrac d{dt}(x^2+y^2)=2\lambda(x^2+y^2)$。特别地，$\lambda=0$ 时长度不变。

\begin{proof}

由乘积法则，导数为 $2xx'+2yy'$。代入速度场，两个交叉项 $-2\omega xy$ 与 $2\omega xy$ 抵消，剩下 $2\lambda(x^2+y^2)$。

因此转动不会单独改变长度，正负径向项则分别使长度平方增长或衰减。这是本线性系统的精确计算，并非任意速度场都可如此分离。

\end{proof}
